\section{Idea 4: decompose by majority projection}\label{sec:components}
Join two orders of \(D\) by an edge when they differ by one adjacent
swap. Let its connected components be \(K\). Connectivity here always
means paths in this adjacent-swap graph.
\lean{\path{Condorcet.AdjacentSwap}, \path{Connected},
\path{component} and \path{components}. These are not definitions
using an arbitrary graph on the domain.}

\begin{lemma}[Majority closure]\label{lem:median}
A full domain contains the majority order of every three of its orders.
\end{lemma}
\begin{proof}
The majority relation is a linear order by the Condorcet property.
It preserves each chosen triple law. A never-middle law equates two
comparison bits, so majority preserves their equality. For a never-first
law, if majority put the distinguished label above both others, two
sets of at least two of the three voters would support the two
comparisons. Their intersection contains a voter putting that label
first, a contradiction. Reverse the argument for never-last. Fullness
then puts the majority order in the family. Fixed pair comparisons,
when present, are preserved too.
\end{proof}
\lean{\path{Condorcet.medianOrder} constructs an actual strict order
from semantic majority; \path{IsFull.median_mem} proves membership.
The Boolean-coordinate statement is \path{IsFull.encoded_medianClosed}.}

\begin{lemma}[A component is determined by its fixed comparisons]\label{lem:projection}
Let \(D\) be full and \(K\) one component. Let \(U\) be the pair
comparisons having the same value throughout \(K\). Then
\[
 K=\{\tau\in D:\tau\text{ agrees with }U\}.
\]
\end{lemma}
\begin{proof}
Only the reverse inclusion needs proof. Suppose \(\tau\in D\) agrees
with \(U\), and choose \(\sigma\in K\) minimizing the number of pair
comparisons on which it disagrees with \(\tau\). If they differ,
some differing comparison is not fixed on \(K\), so choose
\(\rho\in K\) agreeing with \(\tau\) on that comparison. Project
a path in \(K\) from \(\sigma\) to \(\rho\) by
\[
 v\longmapsto\maj(\sigma,\tau,v).
\]
Every projected order lies in \(D\) by Lemma~\ref{lem:median}.
Consecutive projections differ in at most the one comparison changed
by the original edge. Two linear orders differing in one comparison
are adjacent by a swap: if another label lay between the exchanged
labels, transitivity would force some further comparison to change.
Thus the projected path stays in \(K\) and starts at \(\sigma\).
Where \(\sigma,\tau\) agree, projection keeps that value. Where
they disagree, projection follows \(v\). Its final point corrects
the selected disagreement and introduces no disagreement on which
\(\sigma\) was already correct. It is therefore closer to \(\tau\),
a contradiction.
\end{proof}
\lean{\path{Condorcet.IsFull.component_from_fixed} and
\path{component_membership}. The formal proof encodes one Boolean
coordinate per unordered pair; \path{encode_injective} proves that
this representation loses no order information.}

\begin{theorem}[Component reduction and count]\label{thm:components}
If \(D\) is full, every adjacent-swap component \(K\) lies in \(\F_n\),
and
\[
 \boxed{\sum_K|K|=M,\qquad\sum_K\Sig(K)=\Sig(D)+G(D).}
\]
\end{theorem}
\begin{proof}
A chosen never-middle law \(xN_2\) equates the two comparisons of
\(x\) with the other labels of that triple. A swap changes at most
one of those bits, so their common value cannot change along a path.
Throughout a component, \(x\) is consequently above both labels or
below both. Replace this law by \(xN_3\) in the first case and \(xN_1\)
in the second, and impose all the comparisons \(U\) fixed on the
component. Retain the original extreme laws. Every order in \(K\)
satisfies the resulting constraints. Conversely an order satisfying
them satisfies the old laws, since the fixed comparisons recover
each replaced never-middle law. It lies in \(D\) and agrees with
\(U\), hence lies in \(K\) by Lemma~\ref{lem:projection}. This proves
\(K\in\F_n\), including the required fullness and nonemptiness.

Components partition \(D\), giving the first identity. In one original
insertion fibre, a run lies in a single component by its consecutive
swap edges. A component cannot contain two separated original runs:
its own fibres are intervals by Lemma~\ref{lem:interval}. Thus every
run belongs to exactly one component and contributes one deletion
output there. Counting all runs gives \(\sum_K\Sig(K)=\Sig+G\).
\end{proof}
\lean{\path{Condorcet.IsFull.component_inF} proves the single-component
statement; \path{IsFull.components_inF} applies it to members of
\path{components}. The identities are \path{sum_components_card} and
\path{IsFull.sum_components_deletionSum}. Their premises are fullness,
not maximality; maximality is needed later for the gap inequality.}

\begin{example}[What the extra count means]
The full domain \(\{123,132,231,321\}\), defined by \(1N_2\), has
components \(\{123,132\}\) and \(\{231,321\}\). Its deletion sum is
6 and its gap count 2. Each component has deletion sum 4, so the
component sums add to \(8=6+2\). Splitting the family counts the
previously separated fibre runs separately.
\end{example}
